% ==============================================================================
% CHƯƠNG 2: CƠ SỞ LÝ THUYẾT VÀ THIẾT KẾ CHUNG
% ==============================================================================
\chapter{Cơ sở lý thuyết và thiết kế chung}
\label{chap:co-so}

\section{Đồ thị, đường đi và bài toán đường đi ngắn nhất}
\label{sec:do-thi}

\begin{dinhnghia}[Đồ thị có hướng có trọng số]
  Đồ thị có hướng có trọng số là bộ $G = (V, E, c)$ với $V$ là tập đỉnh hữu hạn,
  $E \subseteq V \times V$ là tập cung và $c : E \to \R_{\ge 0}$ gán cho mỗi cung $(u, v)$
  một chi phí không âm $c(u, v)$.
\end{dinhnghia}

Một \thuatngu{đường đi}[path] từ $s$ tới $t$ là dãy đỉnh $P = (v_0, v_1, \dots, v_k)$ với
$v_0 = s$, $v_k = t$ và $(v_{i-1}, v_i) \in E$ với mọi $i$ \cite{rosen2019}. Chi phí của đường
đi là
\begin{equation}
  c(P) = \sum_{i=1}^{k} c(v_{i-1}, v_i).
  \label{eq:chi-phi-duong}
\end{equation}
Bài toán đường đi ngắn nhất tìm $P^\ast$ sao cho $c(P^\ast) = h^\ast(s)$, trong đó
$h^\ast(v)$ là chi phí nhỏ nhất từ $v$ tới đích $t$ (bằng $+\infty$ nếu không tới được).

\section{Thuật toán A*}
\label{sec:astar}

A* duy trì tập biên \OPEN{} các đỉnh đã sinh nhưng chưa mở rộng, và luôn chọn đỉnh có
\begin{equation}
  \NoiBat{f(n) = g(n) + h(n)}
  \label{eq:f}
\end{equation}
nhỏ nhất, trong đó $g(n)$ là chi phí tốt nhất đã biết từ $s$ tới $n$ và $h(n)$ là ước lượng
chi phí còn lại từ $n$ tới đích \cite{hart1968, russell2021}.

\begin{dinhnghia}[Heuristic chấp nhận được và nhất quán]
  Heuristic $h$ là \nhanmanh{chấp nhận được} nếu $h(n) \le h^\ast(n)$ với mọi $n$, và là
  \nhanmanh{nhất quán} nếu $h(t) = 0$ và $h(u) \le c(u, v) + h(v)$ với mọi cung $(u, v) \in E$.
\end{dinhnghia}

\begin{dinhly}[Tính tối ưu của A*]
  \label{dl:toi-uu}
  (i) Mọi heuristic nhất quán đều chấp nhận được. (ii) Nếu $h$ chấp nhận được thì A* (có mở
  lại đỉnh) trả về đường đi tối ưu. (iii) Nếu $h$ nhất quán thì giá trị $f$ của các đỉnh được
  mở rộng không giảm và mỗi đỉnh được mở rộng tối đa một lần.
\end{dinhly}

\begin{loigiai}[Chứng minh (i)]
  Xét đường tối ưu $n = v_0, v_1, \dots, v_k = t$. Cộng các bất đẳng thức
  $h(v_{i-1}) - h(v_i) \le c(v_{i-1}, v_i)$ với $i = 1..k$ ta được
  $h(n) - h(t) \le h^\ast(n)$, mà $h(t) = 0$ nên $h(n) \le h^\ast(n)$. Nếu $n$ không tới được
  đích thì $h^\ast(n) = +\infty$ và bất đẳng thức hiển nhiên đúng.
\end{loigiai}

Phần (ii) và (iii) là kết quả kinh điển của Hart, Nilsson và Raphael \cite{hart1968}; với
$h \equiv 0$, A* trở thành thuật toán Dijkstra.

\section{Các metric không gian}
\label{sec:metric}

Với hiệu tọa độ $\Delta = (\Delta_1, \dots, \Delta_d)$, $d \in \{2, 3\}$, báo cáo dùng ba
metric:
\begin{align}
  h_{M} &= \sum_{i} \abs{\Delta_i}, &
  h_{E} &= \sqrt{\textstyle\sum_{i} \Delta_i^2}, &
  h_{C} &= \max_{i} \abs{\Delta_i}.
  \label{eq:metric}
\end{align}

\begin{dinhly}[Quan hệ giữa ba metric]
  \label{dl:metric}
  Với mọi $\Delta \in \R^d$: $\ h_C \le h_E \le h_M \le d \cdot h_C$ và $h_E \le \sqrt{d}\, h_C$.
\end{dinhly}

\begin{loigiai}
  Đặt $a = \max_i \abs{\Delta_i}$. Ta có $a^2 \le \sum \Delta_i^2 \le d a^2$ nên
  $h_C \le h_E \le \sqrt{d}\, h_C$. Mặt khác
  $\big(\sum \abs{\Delta_i}\big)^2 = \sum \Delta_i^2 + 2\sum_{i<j}\abs{\Delta_i}\abs{\Delta_j}
  \ge \sum \Delta_i^2$ nên $h_E \le h_M$, và $h_M \le d \cdot a$.
\end{loigiai}

\begin{ghichu}[Hệ quả dùng trong các task]
  Nếu chi phí mỗi cung không nhỏ hơn khoảng cách theo một metric giữa hai đầu cung thì, nhờ bất
  đẳng thức tam giác, metric đó là heuristic nhất quán; mọi metric nhỏ hơn nó theo
  Định lý~\ref{dl:metric} cũng nhất quán. Metric lớn hơn có thể đánh giá quá chi phí
  thật: tìm kiếm nhanh hơn nhưng mất bảo đảm tối ưu.
\end{ghichu}

\section{Cấu trúc dữ liệu chung \ma{PathNode}}
\label{sec:pathnode}

Lời giải của mọi task được biểu diễn bằng danh sách liên kết đơn theo đúng đặc tả. Ngoài năm
trường bắt buộc, struct có thêm trường \ma{pSafe} đặt \emph{sau} \ma{next} để lưu xác suất an
toàn tích lũy của Task~3; vị trí này giữ nguyên thứ tự khởi tạo của năm trường bắt buộc nên
tương thích với hàm \ma{printPath} của LMS.

\begin{khungma}{PathNode.h}{C++}
struct PathNode {
    std::string name;
    double f;
    double g;
    double h;
    PathNode* next;
    double pSafe;
};

PathNode* createPathNode(const std::string& name, double g, double h);
void appendPathNode(PathNode*& head, PathNode*& tail, PathNode* node);
int pathLength(const PathNode* head);
const PathNode* pathTail(const PathNode* head);
void freePath(PathNode* head);
\end{khungma}

Hàm \ma{createPathNode} luôn đặt $f = g + h$ nên ba giá trị không thể lệch nhau; nối đuôi dùng
con trỏ \ma{tail} nên dựng danh sách $k$ nút tốn $O(k)$.

\section{Lõi A* dùng chung}
\label{sec:loi}

Cả năm hàm đều dựng một đồ thị tường minh \ma{astar::Graph} gồm danh sách kề và một mảng
$h[v]$, sau đó gọi cùng hàm \ma{astar::search}. Nhờ vậy phần khác nhau giữa các task chỉ còn là
\emph{cách dựng đồ thị}, còn phần tìm kiếm được viết và kiểm thử một lần.

\begin{lstlisting}[style=gia, caption={Giả mã lõi tìm kiếm \ma{astar::search}}, label={lst:search}]
function Search($G$, $s$, Goal)
    $g[v] \gets +\infty$ for each $v$;  $g[s] \gets 0$;  parent$[s] \gets$ nil
    push $(s, g = 0, h[s])$ into OPEN
    while OPEN not empty
        $(u, g_u) \gets$ pop entry with smallest ($f$, then $h$, then insertion order)
        if $g_u > g[u]$ then continue            // stale entry (lazy deletion)
        if $u \in$ Goal then return path rebuilt from parent
        for each arc $(u, v)$ with cost $c$
            if $g[u] + c < g[v]$ then
                $g[v] \gets g[u] + c$;  parent$[v] \gets u$
                push $(v, g[v], h[v])$ into OPEN   // may re-open $v$
    return no path
\end{lstlisting}

Các quyết định thiết kế chính:
\begin{dssao}
  \item \textbf{Hàng đợi ưu tiên.} Môi trường không cho dùng \ma{<queue>}, nên \OPEN{} là đống
        nhị phân trên \ma{std::vector} với \ma{std::push_heap}/\ma{std::pop_heap} của
        \ma{<algorithm>}; mỗi thao tác tốn $O(\log |\OPEN|)$.
  \item \textbf{Phá hòa cố định.} Khi $f$ bằng nhau (sai số $10^{-9}$), ưu tiên $h$ nhỏ hơn
        (gần đích hơn), sau đó ưu tiên phần tử được đưa vào trước. Kết quả vì thế tất định.
  \item \textbf{Xóa lười.} Thay vì giảm khóa, mỗi lần cải thiện $g[v]$ ta đẩy một bản ghi mới;
        bản ghi cũ bị bỏ qua khi lấy ra vì $g_u > g[u]$.
  \item \textbf{Cho phép mở lại.} Không dùng tập đóng: đỉnh đã mở rộng vẫn được đưa lại vào
        \OPEN{} nếu tìm được $g$ tốt hơn. Với heuristic nhất quán điều này không bao giờ xảy ra
        (Định lý~\ref{dl:toi-uu}(iii)); với heuristic chỉ chấp nhận được, nó bảo đảm tính tối ưu
        (Định lý~\ref{dl:toi-uu}(ii)).
  \item \textbf{Đầu vào không hợp lệ.} Chỉ số ngoài phạm vi, \ma{mode} sai, điểm đầu/cuối là
        vật cản hoặc không có đường đi đều trả về \ma{nullptr}. Khi điểm đầu trùng đích, danh
        sách chỉ có một nút với $g = 0$.
\end{dssao}

Với heuristic nhất quán, mỗi đỉnh được mở rộng không quá một lần và mỗi cung được nới lỏng
không quá một lần, nên độ phức tạp là $O\big((|V| + |E|)\log |V|\big)$ thời gian và
$O(|V| + |E|)$ bộ nhớ \cite{cormen2022}.

\section{Phương pháp kiểm chứng}
\label{sec:kiem-chung}

Bộ kiểm thử trong \ma{tests/} gồm 23 ca: các ví dụ tính tay cho từng task, các trường hợp biên
(điểm đầu trùng đích, không có đường, chỉ số và chế độ không hợp lệ) và các phép thử ngẫu nhiên
so sánh chi phí của A* với Dijkstra cài độc lập cho Task~1 (30 mạng), Task~4 với Chebyshev
(40 bản đồ) và Task~5 (30 bảng mạch). Script \ma{check_headers.sh} bảo đảm mã nộp chỉ dùng sáu
thư viện được phép, và GitHub Actions chạy lại toàn bộ ở chuẩn C++11 và C++17 sau mỗi lần đẩy mã.
